\documentclass[11pt]{article}

\usepackage[a4paper,margin=2.5cm]{geometry}
\usepackage[utf8]{inputenc}
\usepackage[T1]{fontenc}
\usepackage{lmodern}
\usepackage{textcomp}
\input{glyphtounicode}
\pdfgentounicode=1

\usepackage{amsmath,amssymb,amsthm,mathtools}
\usepackage{authblk}
\usepackage[final,babel=true,activate={true,nocompatibility},protrusion=true,expansion=false]{microtype}
\usepackage{booktabs}
\usepackage{longtable}
\usepackage{array}
\usepackage{enumitem}
\usepackage{xcolor}
\usepackage{listings}
\usepackage{hyperref}

\hypersetup{
  colorlinks=true,
  linkcolor=blue,
  citecolor=blue,
  urlcolor=blue,
  pdftitle={axiom-explorer: an LLM-assisted harness for cross-search over modern axiomatic seeds},
  pdfauthor={Vera Gomez, Francisco Javier},
  pdfsubject={Methodology, LLM-assisted research, open science},
  pdfkeywords={axiom-explorer, LLM-assisted research, open science, methodology, mathematical discovery, preprint, Zenodo}
}

\theoremstyle{definition}
\newtheorem*{principle}{Principle}
\newtheorem*{stop-rule}{Stop rule}
\theoremstyle{remark}
\newtheorem*{remark}{Remark}

% Lightweight typewriter for commands and code
\lstset{
  basicstyle=\ttfamily\small,
  breaklines=true,
  showstringspaces=false,
  columns=fullflexible,
}

\setlist[itemize]{nosep,leftmargin=*}
\setlist[enumerate]{nosep,leftmargin=*}

\title{axiom-explorer: an LLM-assisted harness for cross-search\\
over modern axiomatic seeds in mathematics, with one\\
case study (condensed classifying anima envelope)}

\author[1]{Francisco Javier Vera G\'omez}
\affil[1]{Independent researcher\\
\href{https://orcid.org/0009-0001-3516-5871}{ORCID: 0009-0001-3516-5871}\\
Repository:
\href{https://github.com/Dredok/axiom-explorer}{github.com/Dredok/axiom-explorer}\\
Version: v1.0-preprint\\
DOI: \href{https://doi.org/10.5281/zenodo.20184068}{10.5281/zenodo.20184068}}

\date{\today}

\begin{document}
\maketitle

\begin{abstract}
We describe \emph{axiom-explorer}, an autonomous LLM-assisted experimental
harness that performs systematic cross-search over modern axiomatic seeds in
mathematics. The harness selects a small set of recent, productive axioms
from distinct branches of mathematics, runs a controlled bibliometric and
formal-state sweep over their pairwise combinations, builds dossiers
integrating literature search, finite-case sanity checks (Z3, SymPy), and
where applicable Lean 4 axiomatic skeletons, and filters surviving
candidates against four explicit relevance tests. The role of the large
language model is constrained by explicit stop rules: the LLM is the
orchestrator and an interpretive filter, not an oracle of mathematical
truth, and the human author retains final approval on every release,
publication, and external communication. We describe the architecture,
workflow, confidence ladder, and stop rules; report a single case study (a
quadruple of seeds whose synthesis is reported as a separate companion
preprint); and discuss what the experiment validates, what it does not
attempt to validate, and the methodological and ethical questions
LLM-assisted mathematical-discovery workflows raise. Code, prompts, raw
data, and per-phase reports of the case study are openly available at the
linked repository. \emph{This preprint is itself the output of the
workflow it describes; see~\S\ref{sec:self-disclosure}.}
\end{abstract}

\tableofcontents

\section{Introduction and contribution}\label{sec:intro}

We address a methodological question that the recent rise of capable
language models makes newly tractable:

\begin{quote}
\emph{Can a curated workflow assist a non-specialist human in surfacing
plausible inter-branch connections in current mathematical research, in
a way whose evidential value is bounded and whose risks are explicit?}
\end{quote}

We do not claim to answer it in the affirmative. We propose a concrete
instrument, \emph{axiom-explorer}, which lets the question be asked
empirically over specific seeds. The instrument has three components:

\begin{enumerate}
\item A \textbf{seeded cross-search procedure} that takes a small set of
``modern productive'' axiomatic frameworks from distinct branches and
explores their pairwise (and higher) combinations through five explicit
phases.
\item A \textbf{confidence ladder} that labels every claim it surfaces by
its epistemic status, from L0 (mechanically verified) to L3 (speculative
guess emerging from the cross-search).
\item A \textbf{governed orchestration layer}, where a large language
model orchestrates the workflow under explicit stop rules and the human
author retains all release decisions.
\end{enumerate}

The contribution of this paper is the methodology and its
operationalisation, not a mathematical theorem. We report one case study
(\S\ref{sec:case}) where the methodology surfaced a candidate
inter-branch cardinality envelope; that case study is the subject of a
separate companion preprint, \emph{A conjectural cardinality envelope for
the condensed classifying anima of spectral $\infty$-topoi: a synthesis
across geometric and analogical instances}~\cite{Vera2026Conjecture}, and
we summarise here only the parts of the case study that illustrate the
methodology.

\paragraph{What this paper is not.} This paper is not a proof that
LLM-assisted methodology produces correct mathematics; it cannot be,
since the test of mathematical correctness happens at the level of
individual claims, not of a workflow. It is also not a comparative
evaluation across LLMs, prompting strategies, or tooling choices; we
describe one configuration that ran end-to-end and produced a
checkable, citable artefact, and we do not vary the configuration here.

\paragraph{Audience.} Readers interested in methodology of open
science, LLM-assisted research, philosophy of mathematics, or
mathematical discovery tools. The case study presupposes informal
familiarity with the language of $\infty$-topoi and condensed
mathematics; specialists in those areas are referred to the companion
preprint for the mathematical claims.

\section{Related work}\label{sec:related}

Several lines of prior work motivate axiom-explorer.

\paragraph{Automated conjecture generation in mathematics.} Fajtlowicz's
Graffiti~\cite{Fajtlowicz1988Graffiti}, a system that produced
non-trivial graph-theoretic conjectures from the late 1980s onwards, and
Lenat's AM and EURISKO~\cite{Lenat1976AM} from the same era, are the
classical references for automated mathematical conjecturing. More
recently, the Ramanujan Machine~\cite{RamanujanMachine2021Nature}
produced thousands of conjectured continued-fraction identities; several
have been subsequently proved. AlphaProof and
AlphaGeometry~\cite{DeepMind2024AlphaGeometry, DeepMind2024AlphaProof}
attack the dual problem of \emph{proving} given theorems with the
assistance of LLMs and formal-prover infrastructure. Our work differs:
we do not aim to prove or to conjecture in a single narrow domain, but
to surface inter-branch synthesis candidates that an expert human can
then prove, refute, or recognise as folklore.

\paragraph{Formalisation projects as backbone.} The Mathlib4
library~\cite{Mathlib2020Community} and major formalisation efforts such
as the Liquid Tensor Experiment~\cite{Scholze2022LTE,
LeanCommunity2022LTEFinal}, the formalisation of perfectoid
spaces~\cite{BuzzardCommelinMassot2020Perfectoid}, and the synthetic
Stone duality work in HoTT/Cubical Agda
literature~\cite{CherubiniCoquandGeerligsMoeneclaey2024SSD,
SyntheticZariskiRepo} provide the formal substrate that axiom-explorer
introspects. Without these libraries, the formal-state inventory of
\S\ref{sec:workflow} would have nothing to inspect.

\paragraph{Open science and preprint infrastructure.}
Zenodo~\cite{ZenodoCERN} provides DOI minting, versioning, and a
concept-DOI mechanism that distinguishes individual versions of a
preprint from the work itself. CITATION.cff~\cite{CITATIONcff} and
.zenodo.json~\cite{ZenodoMetadata} provide
machine-readable metadata standards. These are essential for an
LLM-assisted workflow because the workflow can iterate quickly and must
register each iteration honestly.

\paragraph{Disclosure standards for LLM-assisted work.} As of 2026,
major publishers have settled on a position that LLMs are not
co-authors but their use must be disclosed, typically in
acknowledgments or a methods section~\cite{Nature2023LLMPolicy,
ICMJE2023AIGuidance}. Our workflow internalises this requirement: the
disclosure is a structural section of the preprint, not a footnote.

\section{Architecture and workflow}\label{sec:workflow}

We describe axiom-explorer top-down: first what it does at the level of
a research session (Phases 0--5), then how the LLM orchestrator is
governed, then the supporting tooling.

\subsection{Phases of a research session}

A run of axiom-explorer is organised into five sequential phases. Each
phase produces a stable artefact on disk before the next phase begins.

\begin{description}[style=nextline,leftmargin=1.5em]
\item[Phase 0 --- Bibliometric mapping.]
For a chosen quadruple of axiomatic seeds, the harness queries
arXiv across the six binary combinations using five query strategies
(Q1--Q5): strict canonical-phrase AND, grouped OR/AND with synonyms,
primary $\times$ related fanout, author-pair co-publication restricted to
\texttt{cat:math.*}, and single key-author topic search restricted to
\texttt{cat:math.*}. The output is a per-pair density score and a
classification (\emph{saturated}, \emph{warm}, \emph{warm-bridged},
\emph{sparse}, \emph{frontier-bridged}, \emph{virgin}).

\item[Phase 1 --- Formal-state inventory.]
For each seed, the harness clones the relevant formalisation
repositories (Mathlib4 and adjacent Lean/Agda projects) and counts files,
definitions, theorems, lemmas, and explicit TODOs in source headers.
Per-seed gaps are noted as formalisable targets.

\item[Phase 2 --- Binary-combination dossiers.]
For each pair that survives Phase 0 filtering, the harness builds a
dossier integrating literature (the Phase 0 hits), bridging functors or
constructions reported in those papers, and where applicable a
small-scale numerical or symbolic experiment in Python (NumPy, SymPy) or
finite-model SMT in Z3.

\item[Phase 3 --- Deep dives on surviving candidates.]
A candidate that emerges from a binary dossier (typically a conjectured
inequality, identification, or structural pattern) is promoted to a
dossier of its own: precise statement, evidence trail, falsifiability,
explicit confidence level. Where the candidate is a formalisation-shaped
gap, the harness emits a Lean 4 axiomatic skeleton that type-checks the
axioms involved.

\item[Phase 4 --- Filtering against the four relevance tests.]
Each surviving candidate is tested against four criteria:
\textbf{novelty} (not in the literature found in Phase 0),
\textbf{inter-branch connectivity} (touches at least two of the four
seed branches), \textbf{open-problem tangency} (relates to a stated open
problem or programme), and \textbf{non-triviality} (the conjunction of
the two seeds is genuinely required). Candidates failing two or more
tests are demoted.

\item[Phase 5 --- Synthesis and reporting.]
The surviving candidates, together with the methodology meta-analysis
and any forward open questions identified, are written up. Phase 5
output is a draft manuscript that is then iterated against AI peer
review (see~\S\ref{sec:peer-review-loop}) before becoming a public
preprint.
\end{description}

\subsection{Confidence ladder}\label{sec:ladder}

Every substantive claim that survives Phase 4 is tagged on a four-level
ladder, written explicitly into the dossier and propagated into the
preprint:

\begin{itemize}
\item \textbf{L0 --- Verified mechanically}: directly checked in source
code, formal proof, or finite enumeration (Z3 search, Lean type-check).
\item \textbf{L1 --- Strong}: a standard, named result in the cited
literature, with page or theorem reference.
\item \textbf{L2 --- Plausible}: argued or supported by the cited
evidence; not directly verified by the harness or the author.
\item \textbf{L3 --- Speculative}: a guess emerging from the
cross-search, recorded as a question, not an assertion.
\end{itemize}

The ladder is the single most important downstream-protective device in
the workflow. It is the difference between a preprint that promises
folklore-style certainty and one that hands the reader exactly the level
of trust the workflow earned for each claim.

\subsection{LLM orchestrator and stop rules}\label{sec:stop-rules}

The role of the LLM is bounded by four stop rules, encoded in the
operating prompt of the orchestrator:

\begin{stop-rule}
The orchestrator does not publish to Zenodo or any other archive
automatically. The final ``Publish'' click is the human author's.
\end{stop-rule}

\begin{stop-rule}
The orchestrator does not send emails, social-media posts, or any
external communication. Drafts are produced and committed; the human
author dispatches.
\end{stop-rule}

\begin{stop-rule}
The orchestrator does not claim novelty beyond what the preprint
explicitly claims, which is restricted to ``this formulation does not
appear as such in any single source we found''. Claims of unconditional
correctness, proof, or originality of underlying constructions are
suppressed.
\end{stop-rule}

\begin{stop-rule}
The orchestrator does not impersonate mathematical authority. Author
identity, level of expertise, and the LLM-assisted production method are
declared on the title page or in an early section of every preprint
output.
\end{stop-rule}

These stop rules are checked at code-review boundaries: the harness
refuses to mark a candidate as ``done'' for Phase 4 unless the
confidence ladder is filled in for every claim and the production-method
disclosure has been inserted.

\subsection{Tooling}\label{sec:tooling}

The harness is implemented in Python with a small set of dependencies.
Bibliometric queries use the arXiv API (Atom feed). Formal-state
inventory uses git clone (shallow, blob-filter) plus
\texttt{ripgrep} on Mathlib4 and on the Cubical Agda
synthetic-zariski project~\cite{SyntheticZariskiRepo}. Finite-model
checks use Z3~\cite{Z3deMouraBjorner2008}. Symbolic algebra uses
SymPy~\cite{SymPy2017}. Lean axiomatic skeletons use Lean
4~\cite{Lean42021Moura} with the standard library only (no Mathlib4
dependency, to keep the skeletons self-contained). The publication
pipeline uses \texttt{latexmk} for PDF building and a small Python
script that produces a Zenodo-ready ZIP bundle with a
SHA-256-manifested release.

The full software stack and reproducibility instructions are at the
GitHub repository linked in the title block.

\section{The peer-review loop}\label{sec:peer-review-loop}

A distinctive feature of axiom-explorer is that draft preprints are
iterated against multiple independent AI peer reviews before any human
review or any external publication. Concretely: the Phase 5 draft is
submitted to one or more distinct LLM-based reviewers (different model
families, different prompting); each review produces a structured
critique that the orchestrator must address point by point in a new
version of the draft, with explicit retraction of any v1 claims that
the review proved wrong. The trajectory of revisions is tracked in the
acknowledgments section of the final preprint.

In the case study of \S\ref{sec:case}, four AI peer-review rounds
(v1~$\to$~v1.1, v1.1~$\to$~v1.2, v1.2~$\to$~v1.2.1, v1.2.1~$\to$~v1.2.2)
caught three substantive mathematical errors that the original draft had
made:

\begin{itemize}
\item An invariance failure in the proposed size invariant.
\item A cardinality bound on free profinite groups stated in a form
inconsistent with the geometric attestations being cited.
\item A model-theoretic example (the Lascar group of $\mathrm{ACF}_0$)
misclassified as trivial when it is in fact the absolute Galois group
of $\mathbf{Q}$.
\end{itemize}

A fifth round caught a citation error (an arXiv identifier referring to
a paper on turbulence rather than the intended Campion--Cousins--Ye work
on Lascar groups~\cite{CampionCousinsYe2021}). None of these would have
been caught by a single-shot LLM generation; all of them \emph{were}
caught before any human expert reviewer was approached. The loop is not
a substitute for human expert review (\S\ref{sec:limits}) but it is a
cheap and effective first line of defence.

The methodology takes no position on the eventual ceiling of AI peer
review. In the case study, every successive review surfaced strictly
fewer errors than its predecessor, suggesting either convergence or a
shared LLM blind spot; we cannot distinguish these without a human
review.

\section{Case study summary}\label{sec:case}

This section summarises what axiom-explorer surfaced when seeded with a
specific quadruple. Mathematical statements, attestations, falsifiability
criteria, and the structural argument supporting the candidate are
deferred to the companion preprint~\cite{Vera2026Conjecture}.

\subsection{The seed quadruple}

The four chosen seeds, with the rationale for each, were:

\begin{enumerate}
\item \textbf{A1 --- Univalence Axiom and HoTT.} A modern foundational
axiom system distinct from ZFC, with active formal development in
Cubical Agda and Coq~\cite{UniMath, Agda-cubical}.
\item \textbf{A2 --- Condensed mathematics.} The Clausen--Scholze
programme~\cite{ClausenScholze2019Lectures}, a recent (post-2019)
reformulation of topology that has produced striking results in
algebra, functional analysis, and now \'etale homotopy theory.
\item \textbf{A3 --- Perfectoid spaces.} Scholze's
framework~\cite{Scholze2012Perfectoid}, the source of major recent
progress in $p$-adic Hodge theory and the Langlands programme.
\item \textbf{A4 --- Synthetic Ricci curvature.} The CD$(K,N)$ /
Lott--Sturm--Villani synthetic framework for Ricci curvature in metric
measure spaces, an active programme in geometric analysis.
\end{enumerate}

The choice was made by the human author before any LLM cross-search
began, and was not adjusted during the run. The rationale was
\emph{ex ante} that A1--A3 had recent activity at the Scholze school
and A4 was a deliberate near-stranger to provide a negative control.

\subsection{What the run produced}

\paragraph{Phase 0 classification (after Q5 author-topic search added
in iteration v1.1).} Six binary pairs were classified into three
groups:

\begin{itemize}
\item \textbf{Frontier-bridged}: A1$\times$A2 (Univalence $\times$
Condensed) and A2$\times$A3 (Condensed $\times$ Perfectoid). Both showed
the harness's expected pattern: heavy recent activity across distinct
subcommunities, with cross-pair publication via the Scholze school
(A2$\times$A3) or via Cherubini--Coquand and collaborators in synthetic
Stone duality (A1$\times$A2).
\item \textbf{Sparse}: A1$\times$A3 (Univalence $\times$ Perfectoid).
Plausible bridges exist (synthetic algebraic geometry over HoTT) but
the visible literature is thin.
\item \textbf{Virgin}: A1$\times$A4, A2$\times$A4, A3$\times$A4. The
expected negative control held: synthetic Ricci has essentially no
overlap with the other three seeds at the arXiv level.
\end{itemize}

\paragraph{Phase 1 formal-state inventory.}
A1: extensive Cubical Agda development in
synthetic-zariski~\cite{SyntheticZariskiRepo}, none in Mathlib4.
A2: 34 files in \texttt{Mathlib/Condensed/} including
\texttt{Solid.lean} (the solid module class) and \texttt{Cohomology.lean};
the Liquid Tensor Experiment exists as a completed separate
project~\cite{LeanCommunity2022LTEFinal}.
A3: 4 files in \texttt{Mathlib/RingTheory/Perfectoid/}, including a
2025 \texttt{FontaineTheta.lean} with an explicit TODO bridging to
Bhatt's cotangent-complex approach to prismatic cohomology.
A4: no Mathlib4 formalisation found; a small set of
optimal-transport definitions exists but does not connect to
CD$(K,N)$.

\paragraph{Phase 2--3 deep dive.}
The deep dive ran on the A2$\times$A3 pair and surfaced two recent
arXiv preprints that the strict canonical-phrase Q1 search had missed
and that the Q5 author-topic search caught: Haine
et~al.~\cite{HHLMMW2025}, on the condensed homotopy type of a scheme,
and Haine~\cite{Haine2026classifying}, on classifying anima of
condensed $\infty$-categories of points. The second of these announces
a uniform construction whose specialisations recover the pro-\'etale
fundamental group of a scheme (geometric side, A2$\times$A3) and the
Lascar group of a complete first-order theory (model-theoretic side,
adjacent to A4 in spirit but not in literature).

Cross-reading these two papers together with existing
Bergfalk--Lambie-Hanson--\v{S}aroch and Bannister--Basak work on
condensed set theory~\cite{BLHS2023, BannisterBasak2026} yielded the
case study's candidate observation: that the underlying group at
section level of the fundamental group of the condensed classifying
anima of a spectral $\infty$-topos appears to be bounded by
$2^{\kappa_{\mathcal X}}$ for a suitable size invariant
$\kappa_{\mathcal X}$, with two direct geometric attestations
(\textbf{$\mathbf{P}^1_{\mathbf C}$} and \textbf{$\mathbf{P}^1_{\mathbf Q}$})
and several analogical attestations from model theory and set theory.

\paragraph{Phase 4 filtering.}
The candidate passed the four relevance tests (after the AI peer-review
loop refined the formulation): the formulation does not appear in any
single source we found (novelty), it involves at least two of the four
seed branches (inter-branch connectivity), it relates to the announced
Haine--Damaj--Zhang programme on Lascar groups as classifying anima
(open-problem tangency), and the condensed setting is non-trivially
required (non-triviality).

The candidate, the structural argument supporting it, the falsifiability
section, and the explicit confidence labelling are reported in the
companion preprint~\cite{Vera2026Conjecture} at confidence
\textbf{L2~--~Plausible}.

\subsection{What the run did not produce}

\begin{itemize}
\item No proof. The candidate is a conjecture, registered for
community feedback.
\item No new mathematical \emph{construction}. The classifying-anima
unification it sits on is due to Haine~\cite{Haine2026classifying}; the
geometric attestations are due to Haine
et~al.~\cite{HHLMMW2025}; the Lascar-group identification is due to
Campion--Cousins--Ye~\cite{CampionCousinsYe2021}; the condensed
set-theoretic context is due to several recent
papers~\cite{BLHS2023, BannisterBasak2026, BergfalkLH2024}. The
synthesis contribution is the \emph{formulation of a uniform
cardinality envelope across} those constructions, recorded as a
falsifiable conjecture.
\item No claim about LLM-orchestrated mathematics in general. We
report what happened in one run. The methodology section above
describes the configuration; the case study reports its output. A
generalisation to other seeds, other domains, or other LLM
configurations is future work.
\end{itemize}

\section{Discussion: limits and risks}\label{sec:limits}

We separate methodological limits from ethical and sociological risks.

\subsection{Methodological limits}

\begin{itemize}
\item \textbf{Shared LLM blind spots.} The AI peer-review loop catches
many errors but cannot reliably catch the errors all LLM reviewers share.
In our case, all reviewers were derived from a small set of foundation
models with overlapping pre-training. A single human expert review on
the eventual preprint is necessary to close this gap; our preprint is
explicitly an invitation for that review.
\item \textbf{Selection bias in the bibliometric layer.} The arXiv
sweep finds papers indexed on arXiv. Important work published in
non-arXiv venues, in non-English sources, or whose terminology has
drifted (the harness missed several papers in iteration v1 that were
caught in v1.1 only after the related-terms vocabulary was expanded by
the author) is silently absent. The harness reports this honestly but
cannot fix it.
\item \textbf{Confidence ladder is not calibrated.} L2 means
``plausible to us'', not ``probability $> p$'' for any specific
$p$. We do not have empirical data to calibrate the levels, only an
ordinal scale.
\item \textbf{One case study is not validation.} A single run that
surfaced a single candidate does not prove the methodology generalises.
It proves the methodology terminates with an artefact.
\end{itemize}

\subsection{Ethical and sociological risks}\label{sec:ethics}

\begin{itemize}
\item \textbf{Reputation laundering.} An LLM-assisted workflow could be
used to publish quickly with thin substance, leveraging the author's
identity to lend credibility to LLM output. The four stop rules
(\S\ref{sec:stop-rules}) and the explicit production-method disclosure
in every preprint are the structural defence; we do not claim they are
sufficient.
\item \textbf{Citation-graph pollution.} Issuing many low-quality
preprints with DOIs pollutes the citation graph. The methodology
encourages \emph{registered conjectural observations}, not unlimited
output. We have run a single case study and expect to run at most a
small number more before reassessing.
\item \textbf{Asymmetric burden on human reviewers.} A workflow that
produces preprints faster than humans can review them shifts review
burden to the community. We address this by (a) explicitly stating
that the preprint is intended as a falsifiable observation, not as a
new theorem; (b) keeping the preprint short; (c) listing the four
ways the conjecture could be falsified (folklore, counterexample, scope
sharpening, contextualisation) so a reviewer can engage at low cost.
\item \textbf{LLM authorship boundary.} We follow the emerging
consensus~\cite{Nature2023LLMPolicy, ICMJE2023AIGuidance} that LLMs are
not authors. The orchestrator is a tool whose use is disclosed, with
the same epistemological status as Z3, SymPy, or a literature
database.
\end{itemize}

\section{Production-method disclosure}\label{sec:self-disclosure}

This preprint is itself the output of the workflow it describes. The
text was drafted by an LLM orchestrator (Claude, Anthropic; precise
model version recorded in the repository commit history) operating
under the four stop rules of~\S\ref{sec:stop-rules}, iterated against
AI peer review, and committed to git with the human author's explicit
approval at each merge boundary. The mathematical references are
verified by the human author; the methodology claims (the stop rules
in particular) are committed by the human author. As with the
companion preprint, the human author retains full responsibility for
the artefact.

\section{What is in the repository}\label{sec:repo}

The full reproducibility package is at
\url{https://github.com/Dredok/axiom-explorer}:

\begin{itemize}
\item \texttt{paper/methodology/} --- the source of this preprint.
\item \texttt{paper/conjecture/} --- the source of the companion
preprint~\cite{Vera2026Conjecture}.
\item \texttt{src/axiom\_explorer/} --- Python implementation of the
harness phases.
\item \texttt{research/phase0-bibliometric/} through
\texttt{research/phase5-pi-n-cond/} --- per-phase reports of the
case study, including raw arXiv search dumps and dossiers.
\item \texttt{lean/AxiomExplorer/} --- a Lean~4 axiomatic skeleton of
Synthetic Stone Duality, built clean.
\item \texttt{docs/} --- methodology documents and the
case-study synthesis.
\item \texttt{.mailmap} --- canonical author-identity mapping.
\end{itemize}

\section{Future work}\label{sec:future}

Concrete directions that we leave for v2 of the methodology or for a
second case study:

\begin{itemize}
\item A second case study, with a deliberately different seed
quadruple, to test whether the workflow surfaces different kinds of
candidates from different seeds, or only the same shape of finding.
\item Calibration of the confidence ladder: in particular,
distinguishing L2 (plausible) from L3 (speculative) by an explicit
operational test rather than by orchestrator judgement.
\item Cross-LLM peer review: running the same Phase 5 draft through
reviewers from genuinely independent model families (different
pre-training data, not only different fine-tunings), to test the
shared-blind-spot hypothesis of~\S\ref{sec:limits}.
\item Human expert review pre-publication: requiring an expert reading
before the Zenodo deposit. We did not require this in the case study
because the preprint is framed as an invitation to such review; a
future methodological iteration could move the human review forward in
the pipeline.
\end{itemize}

\section*{Acknowledgments}

We thank the multiple anonymous AI reviewers who critiqued successive
drafts of the case study and of this methodology paper, and whose
specific corrections are credited in detail in the companion
preprint~\cite{Vera2026Conjecture}. We thank the developers of arXiv,
Mathlib4, Zenodo, Lean, Z3, and SymPy, on whose open infrastructure
this workflow is built. The error of any methodological framing in
this paper is the author's.

\bibliographystyle{plain}
\bibliography{references}

\end{document}
